\chapter{Hyperpigmentation}

\section*{Definition}
Hyperpigmentation refers to darkening of the skin due to increased melanin production or deposition, either localized or diffuse. Common patterns include epidermal (sun-related), dermal (drug-induced), and mixed (post-inflammatory) hyperpigmentation.

\section*{Pathophysiology}
Melanin synthesis in melanocytes is catalyzed by tyrosinase and stimulated by \ensuremath{\alpha}-MSH, UV radiation, inflammation, and hormones. Increased melanocyte activity (melanosis) or number (melanocytosis) produces hyperpigmentation. Post-inflammatory hyperpigmentation results from melanin transfer to keratinocytes during inflammation; dermal pigmentation occurs when melanin deposits in macrophages (melanophages) in the dermis, producing a slate-gray hue.

\section*{Etiology}
\begin{itemize}
  \item \textbf{UV-induced:} Sun exposure, tanning beds, lentigines, ephelides (freckles)
  \item \textbf{Hormonal:} Melasma (pregnancy, oral contraceptives), Addison's disease, Nelson syndrome
  \item \textbf{Post-inflammatory:} Acne, eczema, psoriasis, burns, lichen planus, trauma
  \item \textbf{Drug-induced:} NSAIDs, antimalarials, amiodarone, minocycline, phenytoin, chemotherapeutic agents
  \item \textbf{Systemic:} Hemochromatosis (bronze diabetes), primary biliary cholangitis, POEMS syndrome
  \item \textbf{Metabolic:} Vitamin B12 deficiency, pellagra (niacin deficiency), acanthosis nigricans
\end{itemize}

\section*{Indications for Evaluation}
Seek care if:
\begin{itemize}
  \item Diffuse hyperpigmentation with fatigue, weight loss, or hypotension (suspect Addison's disease)
  \item New, changing, or irregularly bordered pigmented lesions (rule out melanoma)
  \item Mucosal or palmar hyperpigmentation suggesting systemic disease
  \item Pigmentation associated with systemic symptoms: weakness, abdominal pain, neuropsychiatric changes
\end{itemize}

\section*{Related Symptoms}
\begin{itemize}
  \item Melasma, lentigines, freckles, post-inflammatory hyperpigmentation, Addisonian pigmentation
\end{itemize}
\textbf{Note:} Wood's lamp examination enhances visualization of epidermal melanin and can help distinguish epidermal from dermal hyperpigmentation.

\section*{Self-Care / Home Management}
\begin{itemize}
  \item Daily broad-spectrum sunscreen (SPF \ensuremath{\geq}50) with UVA/UVB protection; reapply every 2--3 hours
  \item Sun-protective clothing, wide-brimmed hats, and avoidance of peak UV hours (10 AM -- 4 PM)
  \item Topical brightening agents: vitamin C (L-ascorbic acid), kojic acid, azelaic acid 15--20\%
\end{itemize}

\section*{Diagnostic Approach}
\begin{itemize}
  \item History of onset, triggers, medications, family history, and associated systemic symptoms
  \item Wood's lamp examination, dermatoscopy, and skin biopsy if uncertain diagnosis
  \item Laboratory workup: ACTH, cortisol, iron studies, B12, ANA, and thyroid function as guided by presentation
\end{itemize}

\section*{Treatment / Management Overview}
\begin{itemize}
  \item First-line: Hydroquinone 2--4\% cream (limited to 3--4 months to avoid ochronosis), retinoic acid, and sunscreen combination
  \item Second-line: Tretinoin, azelaic acid, kojic acid, and topical corticosteroids for melasma
  \item Procedures: Chemical peels (glycolic, salicylic, TCA), Q-switched laser, microneedling, IPL
  \item Systemic: Treat underlying endocrinopathy, discontinue offending drugs, manage systemic disease
\end{itemize}

\clearpage
